\subsection{REGULARITY OF A COMPACT SPACE}

PROPOSITION I. Let X be a compact space, x a point of X. In order that a filter base B formed of closed neighbourhoods of x should be a fundamental system of neighbourhoods of x it is necessary and sufficient that the intersection of the sets of B consists of x alone.

The condition is necessary since X is Hausdorff (§ 8, no. 1, Proposition 1). It is sufficient, for it signifies that x is the only cluster point of B; hence B converges to x by the Corollary to Theorem 1 of no. 1.

COROLLARY. Every compact space is regular.

For it follows from axiom (Hi) (§ 8, no. 1, Proposition 1) that the filter base formed by all the closed neighbourhoods of an arbitrary point of the space satisfies the condition of Proposition 1.

The following proposition amplifies the Corollary to Proposition 1:

PROPOSITION 2. Let X be a compact space and let A, B be two disjoint closed subsets of X. Then there exist two open sets U, V, such that \(U \cap V = \varnothing\) and \(A \subset U\) and \(B \subset V\) .

Suppose the conclusion is false. If every neighbourhood U of A meets every neighbourhood V of B, then the sets \(U \cap V\) form a filter base B on X, which therefore has a cluster point \(x \in X\) . Now x must lie in A, since if y is any point of X not in A there is a neighbourhood of y and a neighbourhood of A which do not intersect, since X is regular, and thus y cannot be a cluster point of B. Similarly x must lie in B and we have a contradiction.

\pandocbounded{\includegraphics[keepaspectratio]{images/part001_de955d1b.jpg}}

This proposition has important consequences which will be examined in Chapter IX, § 4.

The non-Hausdorff quasi-compact space X of Example 2 of no. 1 does not satisfy axiom (O \(_{III}\) ), nor a fortiori the property stated in Proposition 2, since any two non-empty open sets of this space always intersect.

